\section{Extended Related Work}
\label{app:related}

\paragraph{Static and head-level attention specialization.} Studies on attention dynamics reveal that token utility is structurally skewed across heads and positions. \citet{streamingllm} identify ``attention sinks,'' demonstrating that retaining a few initial tokens allows for stable streaming inference without the full KV cache. Building on this, \citet{fastgen} propose profiling attention heads to statically categorize them for distinct retention policies. Mechanistic studies \cite{retrievalhead} further reveal that specific ``retrieval heads'' are uniquely critical for long-range dependencies. Leveraging these characteristics, \citet{razorattn}, \citet{duoattn}, and \citet{prulong} enhance inference efficiency by statically identifying retrieval-critical heads to retain full context, while restricting others to local context. Unlike these static heuristics, our approach introduces a learnable, input-dependent admission mechanism that captures dynamic token utility patterns at runtime.

\paragraph{Read-time KV Selection.} A large body of work accelerates long-context inference by selectively attending to a sparse, query-dependent subset of cached KV pairs at each step. These methods approximate attention while maintaining the complete underlying cache for future queries. Existing approaches exploit sparsity in various dimensions, including temporal \cite{hip,specache,attnpredictor}, channel \cite{sparq}, inter-layer \cite{omnikv,hshare}, head-wise \cite{moh,qada}, top-$p$ \cite{twilight}, and low-rank \cite{sea,loki} structures. Some retrieval-oriented designs treat the KV cache as an indexed memory \cite{arkvale,retrievalattn,squeezedattn,pqcache,clusterkv,magicpig,retroinfer}. Systems-oriented designs further integrate selection with KV-aware paging \cite{quest,lserve}, runtime management \cite{infinigen,shadowkv}, sparse prefill patterns \cite{minference,flexprefill,spargeattn,xattn}, or learnable indexers \cite{seerattn,moba,nsa,seerattnr,deepseekv32,deepseekv4}. Our write-time admission mechanism is complementary to these read-time strategies, allowing for compound efficiency gains by filtering redundant states prior to selection.

\paragraph{Post-write KV Eviction.} Another line of work focuses on maintaining a bounded cache by retrospectively pruning previously accumulated states. Representative methods employ statistics derived from past attention \cite{scissorhands,h2o,tova} or intrinsic token properties \cite{sirllm,l2norm,qfilters,expectedattn} to evict tokens. Other approaches refine the eviction criteria through learnable interaction scores \cite{dcp,attngate}, importance estimation \cite{locret,kvzap,kvp,trimkv,fastkvzip}, prompt-guided heuristics \cite{snapkv,nacl,pyramidinfer,finch,andpro}, reconstruction objectives \cite{kvzip}, or reasoning trace redundancy \cite{dms,headkv,rkv}. Unlike Admission, which pre-filters redundant states at the source, Eviction manages cache bounds by removing obsolete history; thus, the two serve distinct, complementary roles in the token lifecycle.

\paragraph{Representation-level KV compression.} A parallel line of research seeks to compress the KV representations themselves. Instead of discarding tokens, merging-based methods aggregate information by combining keys and values based on semantic similarity \cite{kvmerger}, weighted accumulation \cite{cam}, low-rank residual recovery \cite{less}, or learnable mechanisms \cite{dmc,deepseekv4}. Additionally, dimensionality reduction techniques exploit inherent redundancy in hidden states by projecting KV pairs into lower-rank subspaces \cite{eigenattn,shadowkv,palu,think,matryoshkakv}, leveraging inter-layer redundancy \cite{minicache,lckv,cla,xkv}, or applying quantization \cite{kvquant}. These lines of work are largely orthogonal to our admission mechanism.